%!TEX root = ./../../thesis-trier.tex

Here, we collect results about transformation semigroups generated by a permutation group~$G$ and a set $T \subseteq \mathcal T_n$ of non-permutations.

\begin{proposition}
	\label{prelim:prop:group_of_units}
	Let $G \le_G \mathcal S_n$ and $T \subseteq \mathcal T_n \setminus S_n$. 
	Then $G$ is the group of units of $\langle G \cup T \rangle_S$.
\end{proposition}
\begin{proof}
		By \crefitem{prelim:prop:rank-defect}{prelim:prop:item:rank-eqn},
		for $t \in \mathcal T_n \setminus \mathcal S_n$
		or $h \in \mathcal T_n$
		we have $\concat{t}{h}, \concat{h}{t}  \in \mathcal T_n \setminus S_n$,
		which gives the claim.
\end{proof}

By successively using the formula $gt = t^g g$ for a transformation $t$
and permutation $g$, we obtain the next statement.

\begin{proposition} 
	\label{sync_max:prop:factorization_KG}
	Let $G \le_G \mathcal S_n$ be a permutation group and $T \subseteq \mathcal T_n$. 
	Set $S = \langle G \cup T \rangle_S$. Let~$K = \{ t^g \mid t \in T, g \in G \}$ be the set of conjugates of elements from~$T$.
	Then $S = \langle K \rangle_M G = G \langle K \rangle_M$.
\end{proposition}
\begin{proof}
	As $K \subseteq S$ and $G \subseteq S$,
	we have $\langle K \rangle_M \subseteq S$
	and so also $\langle K \rangle_M G \subseteq S$.
	Next, suppose $s \in S$. Then we have $s = g_0 t_1 g_1 t_2 g_2 \cdots t_{n-1} g_{n-1} t_n g_n$
	and 
	\begin{equation} \label{sync_max:prop:factorization_KG:eqn}
		g_0 t_1 g_1 t_2 g_2 \cdots t_{n-1} g_{n-1} t_n g_n 
		= t_1^{g_0} t_2^{g_0 g_1} \cdots t_{n-1}^{g_0 g_1 \cdots g_{n-2}} t_n^{g_0 g_1 \cdots g_{n-1}} g_0 g_1 \cdots g_n.
	\end{equation}	
	The above formula can be seen by induction on the number of~$t_i$'s
	in the product. If~$s$ is a product of permutations, then 
	the formula clearly holds.
	Otherwise,
	we have 
	\[ 
	g_0 t_1 g_1 t_2 g_2 \cdots t_{n-1} g_{n-1} t_n g_n 
	= t_1^{g_0} g_0 g_1 t_2 g_2 \cdots t_{n-1} g_{n-1} t_n g_n
	\]
	and, by induction hypothesis,
	\[
	g_0 g_1 t_1 g_2 \cdots t_{n-1} g_{n-1} t_n g_n = t_1^{g_0 g_1} \cdots t_{n-1}^{g_0 g_1 \cdots g_{n-2}} t_n^{g_0 g_1 \cdots g_{n-1}} g_0 g_1 \cdots g_n,
	\]
	which yields \cref{sync_max:prop:factorization_KG:eqn}.
	Since $t_1^{g_0} t_2^{g_0 g_1} \cdots t_{n-1}^{g_0 g_1 \cdots g_{n-2}} t_n^{g_0 g_1 \cdots g_{n-1}} \in \langle K \rangle_M$
	and $g_1 g_2 \cdots g_n \in G$
	we obtain $S \subseteq \langle K \rangle_M G$. So $S = \langle K \rangle_M G$.
	Similarly, we can show
	\[
	 g_0 t_1 g_1 t_2 g_2 \cdots t_{n-1} g_{n-1} t_n g_n 
	 = g_0 g_1 \cdots g_n t_1^{(g_1 g_2 \cdots g_n)^{-1}} t_2^{(g_2 \cdots g_n)^{-1}} \cdots t_{n-1}^{(g_{n-1} g_n)^{-1}} t_n^{g_n^{-1}},
	\]
	which then yields $S = G \langle K \rangle_M$.		
\end{proof}

The following corollary was noted in~\cite[Sec.~2]{DBLP:journals/tcs/AraujoBC13}).

\begin{corollary}
	Suppose $G \le_G \mathcal S_n$ is a permutation group and $t \in \mathcal T_n$.
	If $\langle G \cup \{t\} \rangle_S$
	contains a constant, then some constant %in $\langle G \cup \{t\} \rangle_S$
	is a composition of conjugates of~$t$.
\end{corollary}

Next, we note a special case when $t$ is idempotent.

\begin{corollary}
	\label{sync_max:cor:factorization_EG}
	Suppose $G \le_G \mathcal S_n$ is a permutation group and $t \in \mathcal T_n$ is an idempotent.
	Set $S = \langle G \cup \{t\} \rangle_S$. Let~$E$ be the set of idempotents of $S$.
	Then
	$S = \langle E \rangle_M G$.	
\end{corollary}
\begin{proof}
	Let $K = \{ t^g \mid g \in G\}$.
	By Equation~\eqref{prelim:conj} the mappings in $K$ are idempotents,
	so $K \subseteq E$.
	With Proposition~\ref{sync_max:prop:factorization_KG} 
	we find $S = \langle K \rangle_M G \subseteq \langle E \rangle_M G$.
	As $G \subseteq S$ and $\langle E \rangle_M \subseteq S$,
	and so $\langle E \rangle_M G \subseteq S$,
	we conclude $S = \langle E \rangle_M G$.	
\end{proof} 

As $t_1 \cdots t_n g^{-1} = g^{-1} t_1^g \cdots t_n^g$ for a permutation  $g$ and transformations $t_1, \ldots, t_n$  we have,
in \cref{sync_max:prop:factorization_KG} and \cref{sync_max:cor:factorization_EG},
\[ 
\langle K \rangle_M G = G \langle K \rangle_M \text{ and }
\langle E \rangle_M G = G \langle E \rangle_M.
\]
Further, as $G$ contains the identity transformation, we have
$\langle G \cup \{t\} \rangle_S = \langle G \cup \{t\} \rangle_M$.
However, for a non-permutation $t$, $\langle K \rangle_S G \ne \langle K \rangle_M G$,
as $\operatorname{id}_{[n]} \notin \langle K \rangle_S$.